% ============================================================
% §8b  Rozwiązanie problemu duchów kinetycznych w sektorze solitonowym
% ============================================================
\section{Rozwiązanie problemu duchów kinetycznych
  w~sektorze solitonowym}\label{sec:ghost-resolution}

Sekcja~\ref{sssec:ghost-check} wykazała brak duchów
w~perturbacyjnym sektorze MS-TGP
(stw.~\ref{prop:ghost-free-MS}, stw.~\ref{prop:ghost-free-fundamental}).
Niniejsza sekcja adresuje \textbf{odrębny problem}: osobliwość
kinetyczną profilu solitonowego przy $f(g^*) = 0$, gdzie
$f(g) = 1 + 2\alpha\ln g$ jest funkcją kinetyczną
równania solitonu~\eqref{eq:J-ode}.

\statuslabel{Twierdzenie}: problem duchów jest artefaktem
przybliżenia kontinualnego i~nie występuje
w~pełnym opisie substratowym.

% -------------------------------------------------------------------
\subsection{Sformułowanie problemu}\label{ssec:ghost-problem}
% -------------------------------------------------------------------

\begin{problem}[Osobliwość kinetyczna solitonu TGP]%
\label{prob:kinetic-singularity}
Równanie solitonu TGP (dod.~\ref{app:ogon-masy},
ró{}w.~\eqref{eq:J-ode}):
\begin{equation}\label{eq:ghost-ode}
  f(g)\,g'' + \frac{2}{r}\,g' = V'(g),
  \qquad
  f(g) = 1 + 2\alpha\ln g,
  \quad \alpha = 2,
\end{equation}
posiada \textbf{punkt osobliwy} przy wartości pola
\begin{equation}\label{eq:gstar}
  g^* = e^{-1/(2\alpha)} = e^{-1/4} \approx 0{,}7788,
\end{equation}
w~którym $f(g^*) = 0$. Trzy konsekwencje:
\begin{enumerate}[(i)]
  \item \textbf{Osobliwość ODE}: współczynnik przy $g''$
    zeruje się, równanie degeneruje się
    z~drugiego do pierwszego rzędu.
  \item \textbf{Zmiana znaku}: dla $g < g^*$ zachodzi $f(g) < 0$,
    co odpowiada \textbf{ujemnej energii kinetycznej}
    (,,duch kinetyczny'') w~funkcjonale energii
    $E = \int\!\bigl[\tfrac{1}{2}f(g)(g')^2 + V(g)\bigr]
    4\pi r^2\,dr$.
  \item \textbf{Fizyczna istotność}: profil solitonu mionowego
    i~tauonowego osiąga $g_{\min} < 1$
    (oscylacyjny ogon, ró{}w.~\eqref{eq:J-tail}),
    a~dla dostatecznie dużego $g_0$ profil może
    zbliżyć się do~$g^*$ lub przekroczyć tę wartość.
\end{enumerate}
\end{problem}

\begin{remark}[Zakres problemu]\label{rem:ghost-scope}
Ghost-check MS-TGP (stw.~\ref{prop:ghost-free-MS})
dotyczy jednorodnego tła kosmologicznego $\psi_{\rm bg}(t) > 0$
i~liniowych perturbacji wokół niego. Problem~\ref{prob:kinetic-singularity}
dotyczy \textbf{nieliniowych} konfiguracji solitonowych,
w~których pole $g(r)$ przechodzi przez region $g \approx g^*$.
Są to odmienne reżimy i~wymagają oddzielnego rozwiązania.
\end{remark}

% -------------------------------------------------------------------
\subsection{Rozwiązanie substratowe}\label{ssec:substrate-resolution}
% -------------------------------------------------------------------

Klucz do rozwiązania leży w~rozróżnieniu między
\textbf{pełnym sprzężeniem substratowym} a~jego
\textbf{przybliżeniem kontinualnym}.

\begin{proposition}[Pochodzenie $f(g)$ z~substratu]%
\label{prop:f-from-substrate}
\statuslabel{Twierdzenie}.
Funkcja kinetyczna $f(g) = 1 + 2\alpha\ln g$ jest
rozwinięciem logarytmicznym substratowej funkcji
kinetycznej $K_{\rm sub}(g)$ wokół próżni $g = 1$,
obowiązującym w~reżimie $|g - 1| \ll 1$.

\medskip\noindent
\textbf{Wyprowadzenie.}
W~substracie $\Gamma = (V, E)$ geometryczne sprzężenie
kinetyczne ma postać (tw.~\ref{prop:substrate-action},
ró{}w.~\eqref{eq:K-geometric}):
\begin{equation}\label{eq:Kij-substrate}
  K_{ij} = J\,(\varphi_i\,\varphi_j)^2
  \qquad (\text{link } \langle ij \rangle \in E),
\end{equation}
gdzie $\varphi_i = |\hat{s}_i|$ jest amplitudą parametru porządku
w~węźle~$i$. W~granicy ciągłej, dla jednorodnego tła $\varphi \approx
\varphi_0$, identyfikujemy (aksjomat~\ref{ax:zrodlo}):
\begin{equation}\label{eq:phi-to-g}
  \varphi^2 = \frac{\Phi}{\PhiZero} = g,
  \qquad \varphi = \sqrt{g},
\end{equation}
skąd sprzężenie substratowe w~zmiennej $g$:
\begin{equation}\label{eq:Ksub-g}
  \boxed{
    K_{\rm sub}(g) = K_{\rm geo}\,g^2,
  }
\end{equation}
z~$K_{\rm geo} > 0$ wynikającym z~stałej sprężystości sieci.
\end{proposition}

\begin{remark}[Relacja $K_{\rm sub}(g)$ do $K(\varphi)$]%
\label{rem:K-relation}
W~zmiennych $\varphi$ (sekcja~\ref{ssec:architektura},
tw.~\ref{prop:substrate-action}):
$K(\varphi) = K_{\rm geo}\,\varphi^4$.
Przejście do zmiennej $g = \varphi^2$
daje $d\varphi = dg/(2\sqrt{g})$, więc
$K(\varphi)(\nabla\varphi)^2 = K_{\rm geo}\varphi^4
  \cdot (\nabla g)^2/(4g) = (K_{\rm geo}/4)\,g\,(\nabla g)^2$.
Definiujemy efektywną funkcję kinetyczną
w~zmiennej $g$ jako $K_{\rm sub}(g) = K_{\rm geo}\,g^2$
(z~dokładnością do bezwymiarowej stałej, którą absorbujemy
w~$K_{\rm geo}$). Istotna jest \textbf{struktura potęgowa}:
$K_{\rm sub}(g) \propto g^n$ z~$n = 2 > 0$,
co gwarantuje dodatniość dla $g > 0$.

\smallskip\noindent
\textbf{Uwaga o~lem.~\ref{lem:K_phi2}}: Lemat ten wyprowadza
$K(\varphi) = J a^2 \varphi^2$ z~rozwinięcia gradientowego
hamiltonianu $H_\Gamma = -J\sum(\varphi_i\varphi_j)^2$.
Wynik $\varphi^2$ (nie $\varphi^4$) wynika z~faktu, że
lemat ekstrahuje gradient z~\textbf{energii całkowitej}
$-(\varphi_i\varphi_j)^2$, podczas gdy pełne geometryczne
sprzężenie kinetyczne (tw.~\ref{prop:substrate-action})
definiuje $K_{ij}$ jako wagę członu gradientowego
$K_{ij}(\varphi_i - \varphi_j)^2$, co daje $K = \varphi^4$.
Weryfikacja numeryczna (\texttt{ex163}, sesja v45) potwierdza:
oba opisy s\k{a} poprawne w~swoich ramach ---
lem.~\ref{lem:K_phi2} to \emph{linearyzacja} wok\'o\l{}
jednorodnego t\l{}a $\bar\varphi$,
tw.~\ref{prop:substrate-action} to pe\l{}na \emph{nieliniowa} definicja.
$H_\Gamma \neq F_{\rm kin}$ --- to r\'o\.zne funkcjona\l{}y.
Nie ma sprzeczno\'sci (R11 rozwi\k{a}zany).

Dodatkowy wynik (\texttt{ex163}): ODE kanoniczne ($\alpha = 2$,
$K \propto g^4$) ma barier\k{e} duchow\k{a} $g_{\rm ghost} \approx 0{,}717$,
kt\'ora uniemo\.zliwia selekcj\k{e} $\varphi$-FP
($g_0^\mu = \varphi\,g_0^e > g_{\rm ghost}$).
ODE substratowe ($\alpha = 1$, $K \propto g^2$) nie ma bariery
i~daje $r_{21} = 206{,}77$ z~predykcj\k{a} $m_\tau$ z~precyzj\k{a} $83$~ppm
(rem.~\ref{rem:J-precision-ex157}).
\end{remark}

\begin{remark}[$f(g)$ jako rozwinięcie $K_{\rm sub}$]%
\label{rem:f-as-expansion}
Logarytmiczne rozwinięcie $K_{\rm sub}(g)$ wokół $g = 1$:
\begin{equation}\label{eq:Ksub-expansion}
  K_{\rm sub}(g) = K_{\rm geo}\,g^2
  = K_{\rm geo}\,e^{2\ln g}
  \approx K_{\rm geo}\bigl[1 + 2\ln g + 2(\ln g)^2 + \cdots\bigr].
\end{equation}
Obcięcie do pierwszego rzędu w~$\ln g$
z~normalizacją $K_{\rm geo} = 1$ daje:
\begin{equation}\label{eq:f-truncation}
  K_{\rm sub}(g)\big|_{\rm trunc}
  = 1 + 2\ln g
  \;\;\xrightarrow{\alpha=2}\;\;
  f(g) = 1 + 2\alpha\ln g = 1 + 4\ln g.
\end{equation}
Przybliżenie jest wiarygodne dla $|\ln g| \ll 1$,
tj.\ $g \approx 1$ (reżim próżniowy).
Dla $g \lesssim g^* \approx 0{,}78$
obcięcie nie jest poprawne --- pełne $K_{\rm sub}(g) = g^2$
pozostaje dodatnie.
\end{remark}

% -------------------------------------------------------------------
\subsection{Twierdzenie o~braku duchów w~teorii substratowej}%
\label{ssec:ghost-free-substrate}
% -------------------------------------------------------------------

\begin{theorem}[Brak duchów w~sektorze solitonowym substratu]%
\label{thm:ghost-free-soliton}
\statuslabel{Twierdzenie}.
Niech $H_\Gamma$ będzie hamiltonianem substratu
z~geometrycznym sprzężeniem $K_{ij} = J(\varphi_i\varphi_j)^2 \geq 0$.
Wówczas:
\begin{enumerate}[(i)]
  \item Energia kinetyczna solitonu z~\textbf{pełnym}
    sprzężeniem substratowym:
    \begin{equation}\label{eq:T-substrate}
      T = \frac{1}{2}\,K_{\rm sub}(g)\,(\nabla g)^2
        = \frac{K_{\rm geo}}{2}\,g^2\,(\nabla g)^2
        \geq 0
        \qquad \forall\; g > 0.
    \end{equation}
  \item Równanie solitonu z~$K_{\rm sub}(g)$ nie posiada
    osobliwości kinetycznej:
    \begin{equation}\label{eq:soliton-substrate}
      \boxed{
        K_{\rm geo}\,g^2\,g''
        + K_{\rm geo}\,g\,(g')^2
        + \frac{2}{r}\,K_{\rm geo}\,g^2\,g'
        = V'(g).
      }
    \end{equation}
    Współczynnik przy $g''$ wynosi $K_{\rm geo}\,g^2 > 0$
    dla każdego $g > 0$ --- \textbf{brak degeneracji}.
  \item Żadna konfiguracja solitonowa $g(r) > 0$
    nie generuje ujemnej energii kinetycznej.
\end{enumerate}
\end{theorem}

\begin{proof}
\emph{Punkt (i).}
$K_{\rm sub}(g) = K_{\rm geo}\,g^2$; $K_{\rm geo} > 0$
(stała sprężystości sieci); $g > 0$ (fizycznie: $\Phi > 0$
w~sektorze $\mathcal{S}_1$); $(\nabla g)^2 \geq 0$.
Iloczyn trzech czynników nieujemnych jest nieujemny.

\medskip
\emph{Punkt (ii).}
Funkcjonał energii z~$K_{\rm sub}$:
\begin{equation}
  E[g] = \int_0^\infty\!
    \left[\frac{K_{\rm geo}}{2}\,g^2\,(g')^2 + V(g)\right]
    4\pi r^2\,dr.
\end{equation}
Równanie Eulera--Lagrange'a (sferyczna symetria):
\begin{align}
  0 &= \frac{\delta E}{\delta g}
     = K_{\rm geo}\,g\,(g')^2 + V'(g)
     - \frac{1}{r^2}\frac{d}{dr}\!\left[K_{\rm geo}\,g^2\,g'\,r^2\right] \\
    &= K_{\rm geo}\,g\,(g')^2 + V'(g)
     - K_{\rm geo}\bigl[2g(g')^2 + g^2 g'' + \tfrac{2}{r}g^2 g'\bigr].
\end{align}
Upraszczając:
\[
  K_{\rm geo}\,g^2\,g''
  + K_{\rm geo}\,g\,(g')^2
  + \frac{2}{r}\,K_{\rm geo}\,g^2\,g'
  = V'(g).
\]
Współczynnik przy $g''$: $K_{\rm geo}\,g^2 > 0$ dla $g > 0$.
Równanie jest regularne na całym $g \in (0,\infty)$.

\medskip
\emph{Punkt (iii).}
Wynika bezpośrednio z~(i): $T \geq 0$ dla każdego
$g : [0,\infty) \to (0,\infty)$.
Hamiltonian $H = T + \int V(g)\,4\pi r^2\,dr$
ma nieujemny człon kinetyczny, więc nie ma niestabilności
duchowej (Ostrogradski).
\end{proof}

\begin{corollary}[Osobliwość $f(g^*) = 0$ jest artefaktem]\label{cor:ghost-artifact}
Punkt $g^* = e^{-1/4}$, w~którym $f(g^*) = 0$,
nie jest fizyczną osobliwością teorii.
Jest miejscem, w~którym przybliżenie
$K_{\rm sub}(g) \approx 1 + 4\ln g$ traci ważność.
Pełna funkcja substratowa $K_{\rm sub}(g) = K_{\rm geo}\,g^2$
jest regularna i~dodatnia w~tym punkcie:
\begin{equation}
  K_{\rm sub}(g^*) = K_{\rm geo}\,(e^{-1/4})^2
  = K_{\rm geo}\,e^{-1/2} \approx 0{,}607\,K_{\rm geo} > 0.
\end{equation}
\end{corollary}

% -------------------------------------------------------------------
\subsection{Dopasowanie reżimu próżniowego}\label{ssec:matching}
% -------------------------------------------------------------------

Aby obie postacie --- substratowa $K_{\rm sub}(g) = K_{\rm geo}\,g^2$
i~kontinualna $f(g) = 1 + 2\alpha\ln g$ --- dawały identyczne
wyniki w~reżimie próżniowym ($g \approx 1$), wymagamy zgodności
rozwinięć Taylora wokół $g = 1$.

\begin{proposition}[Dopasowanie substrat--kontinuum]%
\label{prop:matching}
\statuslabel{Twierdzenie}.
Niech $\delta = g - 1$ oraz $|\delta| \ll 1$.
\begin{enumerate}[(a)]
  \item \textbf{Substrat}:
    \begin{equation}\label{eq:Ksub-taylor}
      K_{\rm sub}(g) = K_{\rm geo}\,g^2
      = K_{\rm geo}\bigl[1 + 2\delta + \delta^2\bigr].
    \end{equation}
  \item \textbf{Kontinuum}:
    \begin{equation}\label{eq:f-taylor}
      f(g) = 1 + 2\alpha\ln(1+\delta)
      = 1 + 2\alpha\,\delta - \alpha\,\delta^2 + O(\delta^3).
    \end{equation}
  \item \textbf{Dopasowanie do rzędu $\delta^0$}: $K_{\rm geo} = 1$.
  \item \textbf{Dopasowanie do rzędu $\delta^1$}:
    $2K_{\rm geo} = 2\alpha$, skąd $\alpha_{\rm eff} = K_{\rm geo} = 1$.
\end{enumerate}
\end{proposition}

\begin{remark}[Rozbieżność na wyższych rzędach]\label{rem:discrepancy}
Przy rzędzie $\delta^2$: substrat daje $+K_{\rm geo}\,\delta^2 = +\delta^2$,
kontinuum daje $-\alpha\,\delta^2 = -\delta^2$.
Rozbieżność wynosi $2\delta^2$ i~jest \textbf{nieistotna}
w~reżimie próżniowym ($\delta \ll 1$), ale rośnie dla $g \ll 1$ lub $g \gg 1$.
Dokładnie w~tych reżimach opis substratowy staje się konieczny.

Wartość $\alpha = 2$ wyznaczona w~tw.~\ref{thm:alpha-unique}
wynika z~warunku $K(\varphi) = K_{\rm geo}\,\varphi^4$
(pełna zmienna $\varphi$, nie $g$).
Stała $\alpha_{\rm eff} = 1$ w~powyższym rozwinięciu
dotyczy \textbf{innej} parametryzacji --- rozwinięcia
w~$\delta = g - 1$ zamiast w~$\ln g$.
Rozwinięcia te pokrywają się do pierwszego rzędu
i~rozchodzą na wyższych.
\end{remark}

% -------------------------------------------------------------------
\subsection{Konsekwencje fizyczne}\label{ssec:physical-consequences}
% -------------------------------------------------------------------

\begin{remark}[Solitony ciężkie sondują substrat]%
\label{rem:heavy-solitons-substrate}
Profil solitonu z~dużym $g_0 \gg 1$ (ciężkie cząstki: mion, tauon)
osiąga w~regionie oscylacyjnego ogona wartości
$g_{\min} < 1$ bliskie $g^*$ (dod.~\ref{app:ogon-masy}).
W~opisie kontinualnym $f(g) = 1 + 4\ln g$ profil napotyka
osobliwość --- dynamika staje się nieciągła.

W~opisie substratowym:
\begin{enumerate}[(i)]
  \item Soliton z~dowolnym $g_0 > 0$ jest opisany
    regularnym równaniem~\eqref{eq:soliton-substrate}.
  \item Profil $g(r)$ może przechodzić przez $g^*$
    bez osobliwości: $K_{\rm sub}(g^*) = K_{\rm geo}\,e^{-1/2} > 0$.
  \item Cięższa cząstka ($g_0 \gg 1$) głębiej sonduje
    strukturę substratu --- tam, gdzie pełne $K_{\rm sub}(g) = K_{\rm geo}\,g^2$
    istotnie odbiega od przybliżonego $f(g)$.
\end{enumerate}

Hierarchia mas: elektron ($g_0 \approx 1{,}24$) pozostaje
blisko próżni i~jest dobrze opisany przez $f(g)$.
Mion ($g_0 \approx 2{,}0$) i~tauon ($g_0 \approx 2{,}34$)
sondują głębszy reżim. Ich ogon oscylacyjny ($g \approx 1$,
ró{}w.~\eqref{eq:J-tail}) pozostaje \textbf{w~pełni w~reżimie
próżniowym}, gdzie $f(g) \approx K_{\rm sub}(g)$.
\end{remark}

\begin{proposition}[Zachowanie mechanizmu $A_{\rm tail}$]%
\label{prop:Atail-preserved}
\statuslabel{Twierdzenie}.
Mechanizm hierarchii mas oparty na amplitudzie ogona
$A_{\rm tail}(g_0)$ (hip.~\ref{hyp:J-mass-Atail4},
dod.~\ref{app:ogon-masy}) jest \textbf{niezmienniczy}
wobec przejścia $f(g) \to K_{\rm sub}(g)$
w~regionie ogona, ponieważ:
\begin{enumerate}[(i)]
  \item Ogon oscylacyjny $g(r) - 1 \sim B\cos r/r$
    leży w~reżimie $|g - 1| \ll 1$, gdzie
    $f(g)$ i~$K_{\rm sub}(g)$ różnią się o~$O(\delta^2)$.
  \item Stosunek mas $r_{21} = m_\mu/m_e
    = (A_{\rm tail}^{(\mu)}/A_{\rm tail}^{(e)})^4$
    zależy od \textbf{ilorazu} amplitud,
    w~którym korekcje substratowe $O(\delta^2)$
    znoszą się do wiodącego rzędu.
  \item Wynik numeryczny $r_{21} \approx 206{,}77 \pm 1{,}6\%$
    (dod.~\ref{app:ogon-masy}) jest stabilny wobec
    regularyzacji substratowej, ponieważ
    ogony obu solitonów leżą w~tym samym reżimie.
\end{enumerate}
\end{proposition}

\begin{remark}[Podsumowanie statusu ghost-freedom TGP]%
\label{rem:ghost-summary}
Brak duchów w~TGP jest teraz wykazany na \textbf{trzech}
poziomach:
\begin{center}\small
\begin{tabular}{@{}llll@{}}
\toprule
\textbf{Sektor} & \textbf{Wynik} & \textbf{Referencja} & \textbf{Status} \\
\midrule
Zmienne fundamentalne $\psi$
  & $\psi^4(\nabla\psi)^2 > 0$
  & stw.~\ref{prop:ghost-free-fundamental}
  & \statuslabel{Twierdzenie} \\[3pt]
Perturbacje MS-TGP
  & $Q_s = \psi_{\rm bg}^4 > 0$
  & stw.~\ref{prop:ghost-free-MS}
  & \statuslabel{Propozycja} \\[3pt]
Solitony (substrat)
  & $K_{\rm sub}(g) = K_{\rm geo}\,g^2 > 0$
  & tw.~\ref{thm:ghost-free-soliton}
  & \statuslabel{Twierdzenie} \\
\bottomrule
\end{tabular}
\end{center}

\medskip\noindent
We wszystkich trzech sektorach pozytywność energii
kinetycznej wynika z~\textbf{tej samej} przyczyny:
substratowe sprzężenie $K_{ij} = J(\varphi_i\varphi_j)^2$
jest z~definicji nieujemne. Różne sektory
(jednorodne tło, perturbacje, solitony) eksponują
tę samą właściwość w~różnych zmiennych
i~różnych reżimach.
\end{remark}

% -------------------------------------------------------------------
\subsection{Rozstrzygnięcie $\alpha = 1$ vs $\alpha = 2$}%
\label{ssec:alpha-resolution}
% -------------------------------------------------------------------

\begin{result}[Obserwable s\l{}abopolowe s\k{a} niezale\.zne od~$\alpha$]%
\label{res:alpha-independent}
Metryka emergentna $g_{\mu\nu} = (\Phi/\PhiZero)^{2/3}\,\eta_{\mu\nu}$
jest \textbf{aksjomatyczna} (A3): pot\k{e}ga $2/3$ wynika z~geometrycznej
relacji $ds \propto \Phi^{1/3}$ (obj\k{e}to\'s\'c $\propto \Phi$ w~3D),
\textbf{nie} z~dynamiki pola.

W~s\l{}abym polu ($\Phi = \PhiZero(1 + \delta)$, $|\delta| \ll 1$)
operator kinetyczny TGP:
\begin{equation}\label{eq:D-alpha-linear}
  D(\alpha)[\Phi] \;=\; \nabla^2\Phi + \frac{\alpha}{\Phi}(\nabla\Phi)^2
  \;\;\xrightarrow{\;\delta \ll 1\;\;}\;\;
  \PhiZero\,\nabla^2\delta + O(\delta^2).
\end{equation}
Cz\l{}on $(\alpha/\Phi)(\nabla\Phi)^2 = \alpha\,\PhiZero\,(\nabla\delta)^2
+ O(\delta^3)$ jest \textbf{drugiego rz\k{e}du} w~$\delta$.
St\k{a}d linearyzacja jest \textbf{uniwersalna} (niezale\.zna od~$\alpha$):

\medskip
\begin{center}\small
\begin{tabular}{@{}ll@{}}
\toprule
\textbf{Obserwabla} & \textbf{Zale\.zno\'s\'c od $\alpha$} \\
\midrule
$\kappa = 3/(4\PhiZero)$ & niezale\.zna (1.~rz\k{a}d PPN) \\
PPN: $\gamma = 1$, $\beta = 1$ & niezale\.zne (1.~rz\k{a}d) \\
$N_e = \tfrac{1}{3}\ln(1/\varepsilon_0)$ & niezale\.zny (geometryczny) \\
$n_s = 1 - 2/N_e$ & niezale\.zny \\
$r = 12/N_e^2$ & niezale\.zny \\
$K_{\rm Koide}$ & niezale\.zny (algebraiczny) \\
$G_{\rm SM}$ emergence & niezale\.zna (topologiczna) \\
\midrule
ODE solitonu & \textbf{zale\.zna} (re\.zim nieliniowy) \\
\bottomrule
\end{tabular}
\end{center}
\end{result}

\begin{corollary}[$\alpha = 1$ preferowane bez konfliktu]%
\label{cor:alpha1-preferred}
Jedynym sektorem zale\.znym od~$\alpha$ jest ODE solitonu
(re\.zim nieliniowy, $g$ daleko od~$1$).
ODE substratowe ($\alpha = 1$, $K_{\rm sub} = g^2$) jest preferowane
empirycznie: selekcja $\varphi$-FP daje $r_{21} = 206{,}77$
i~$m_\tau$ z~precyzj\k{a} $83$~ppm (rem.~\ref{rem:J-precision-ex157}).
ODE kanoniczne ($\alpha = 2$, $K = g^4$) posiada barier\k{e}
duchow\k{a} $g_{\rm ghost} \approx 0{,}717$
blokuj\k{a}c\k{a} $\varphi$-FP.

Poniewa\.z \.zadna obserwabla s\l{}abopolowa nie zale\.zy od~$\alpha$
(wynik~\ref{res:alpha-independent}), wyb\'or $\alpha = 1$
\textbf{nie koliduje} z~predykcjami grawitacyjnymi ani kosmologicznymi.
Dualizm $\alpha_{\rm grav} \neq \alpha_{\rm sol}$ jest
\textbf{zb\k{e}dny}.
\end{corollary}

\noindent
Weryfikacja numeryczna: \texttt{ex166\_alpha1\_vs\_alpha2.py}
(systematyczne por\'ownanie),
\texttt{ex167\_ppn\_alpha1.py} (linearyzacja PPN).
